\documentclass[11pt,a4paper]{article}

% ---------------------------------------------------------
% PREAMBLE (stable, warning-free, Overleaf-safe)
% ---------------------------------------------------------
\usepackage[utf8]{inputenc}
\usepackage[T1]{fontenc}
\usepackage{lmodern}
\usepackage{amsmath, amssymb, amsfonts}
\usepackage{bm}
\usepackage{geometry}
\usepackage{graphicx}
\usepackage{hyperref}
\usepackage{cite}
\usepackage{authblk}
\usepackage{physics}
\usepackage{mathtools}

\geometry{margin=2.4cm}

\title{\textbf{Companion LVII: Weak-Lensing Tomography in the Relaxation-Driven Cyclic Cosmology}}
\author{Michael Lehmann \\ \texttt{mi.lehmann@gmx.de}}
\date{April 2026}

\begin{document}
\maketitle

\begin{abstract}
Weak-lensing tomography probes the redshift evolution of the gravitational 
potential by correlating shear fields in multiple source bins. In the standard 
cosmological model, the tomographic kernels are determined by the matter power 
spectrum, the growth rate, and the geometry of the Universe. In the 
Relaxation-Driven Cyclic Cosmology (RDCC), the scale-dependent evolution of the 
gravitational potential modifies the lensing kernels, the growth of structure, and 
the redshift dependence of the shear correlations in a characteristic way. This 
Companion develops a continuous description of weak-lensing tomography in RDCC, 
deriving how the relaxation scale influences the tomographic kernels, the 
cross-bin correlations, and the observational signatures in cosmic shear surveys. 
The resulting framework provides a distinctive probe of the RDCC gravitational 
field across cosmic time.
\end{abstract}
% ---------------------------------------------------------
\section{Introduction}

Weak gravitational lensing probes the integrated gravitational potential along the 
line of sight. Tomographic lensing extends this by dividing source galaxies into 
redshift bins, allowing the reconstruction of the redshift evolution of the 
gravitational field. The tomographic kernels encode the geometry of the Universe, 
the matter distribution, and the growth of structure.

In the standard cosmological model, the tomographic kernels evolve smoothly with 
redshift. In RDCC, the gravitational potential evolves differently. The relaxation 
mechanism suppresses the decay of small-scale modes and enhances the coherence of 
large-scale modes. This modifies the tomographic kernels, the growth rate, and the 
cross-bin correlations in a scale-dependent manner, producing a distinctive 
signature in weak-lensing tomography.
% ---------------------------------------------------------
\section{The RDCC Lensing Kernel}

The standard lensing kernel for a source bin $i$ is
\begin{equation}
    W_{i}(\chi)
    = \frac{3H_{0}^{2}\Omega_{m}}{2c^{2}}
      \frac{\chi}{a(\chi)}
      \int_{\chi}^{\chi_{H}} d\chi'\,
      n_{i}(\chi')\,
      \frac{\chi' - \chi}{\chi'}.
\end{equation}

In RDCC, the gravitational potential evolves as
\begin{equation}
    \Phi_{\mathrm{RDCC}}(k,a)
    = \Phi(k,a)
      \left[
        1 - \chi_{\Phi}
        \left(\frac{k}{k_{\mathrm{rel}}}\right)^{\alpha}
      \right].
\end{equation}

The RDCC lensing kernel becomes
\begin{equation}
    W_{i,\mathrm{RDCC}}(\chi)
    = W_{i}(\chi)
      \left[
        1 - \chi_{W}
        \left(\frac{k}{k_{\mathrm{rel}}}\right)^{\alpha}
      \right].
\end{equation}
% ---------------------------------------------------------
\section{Tomographic Power Spectra in RDCC}

The tomographic shear power spectrum between bins $i$ and $j$ is
\begin{equation}
    C_{\ell}^{ij}
    = \int_{0}^{\chi_{H}}
      d\chi\,
      \frac{W_{i}(\chi)\, W_{j}(\chi)}{\chi^{2}}\,
      P_{m}\left(k=\frac{\ell}{\chi},z(\chi)\right).
\end{equation}

In RDCC, the matter power spectrum becomes
\begin{equation}
    P_{m,\mathrm{RDCC}}(k,z)
    = P_{m}(k,z)
      \left[
        1 - \chi_{P}
        \left(\frac{k}{k_{\mathrm{rel}}}\right)^{\alpha}
      \right].
\end{equation}

The tomographic spectrum becomes
\begin{equation}
    C_{\ell,\mathrm{RDCC}}^{ij}
    = C_{\ell}^{ij}
      \left[
        1 - \chi_{C}
        \left(\frac{\ell}{\ell_{\mathrm{rel}}}\right)^{\alpha}
      \right].
\end{equation}
% ---------------------------------------------------------
\section{Cross-Bin Correlations and RDCC Signatures}

The relaxation mechanism enhances the coherence of large-scale modes. The 
cross-bin correlations therefore exhibit:

\begin{itemize}
    \item enhanced correlations at large angular scales,
    \item suppressed correlations at small scales,
    \item a characteristic turnover near $\ell_{\mathrm{rel}}$,
    \item a modified redshift dependence of the cross-bin amplitude.
\end{itemize}

These signatures provide a direct observational probe of the RDCC gravitational 
field across cosmic time.
% ---------------------------------------------------------
\section{Redshift Evolution of the RDCC Kernel}

The redshift evolution of the RDCC kernel is given by
\begin{equation}
    \frac{\partial W_{i,\mathrm{RDCC}}}{\partial z}
    = \frac{\partial W_{i}}{\partial z}
      \left[
        1 - \chi_{W}
        \left(\frac{k}{k_{\mathrm{rel}}}\right)^{\alpha}
      \right]
      - W_{i}\,
        \chi_{W}\,\alpha\,
        \left(\frac{k}{k_{\mathrm{rel}}}\right)^{\alpha}
        \frac{\partial \ln k_{\mathrm{rel}}}{\partial z}.
\end{equation}

This term is unique to RDCC and provides a direct probe of the relaxation scale.
% ---------------------------------------------------------
\section{Conclusion}

Weak-lensing tomography in the Relaxation-Driven Cyclic Cosmology reflects the 
scale-dependent evolution of the gravitational potential and the matter 
distribution. The relaxation mechanism modifies the tomographic kernels, the 
growth rate, and the cross-bin correlations in a scale-dependent manner, producing 
distinctive signatures in cosmic shear surveys. These effects provide a direct 
observational window into the relaxation scale and the temporal structure of the 
RDCC gravitational field. The present Companion establishes the theoretical 
foundation for future studies of cosmic shear, matter evolution, and the interplay 
between light and gravity in RDCC.

\bibliographystyle{apsrev4-2}
\nocite{*}
\bibliography{references}

\end{document}
